% ============================================================
%  Teaching Philosophy - Ronald J. Tackett, Ph.D.
%  Engine: LuaLaTeX
%
%  CHANGELOG
%  v1.0  2026-09-03  Rebuilt in LaTeX from the March 2026 PDF.
%                    Content unchanged; footer date moved to
%                    September 3, 2026 and the closing block from
%                    Spring 2026 to Fall 2026.
% ============================================================
\documentclass[10pt]{article}
\usepackage{tackettdoc}

\setdoclabel{Teaching Philosophy}
\setupdated{September 3, 2026}

\begin{document}

\statementheader
  {Founding Head, School of Foundational Studies}
  {Associate Professor of Physics \textperiodcentered\ Kettering University}
  {Teaching Philosophy}

\lede{Teaching, at its best, is an act of productive disruption. The goal is not to make
things easier for students; it is to make difficulty worthwhile.}

\stsection{I. Where I Start}

I teach physics and foundational science to engineering students, most of whom arrive at
Kettering with strong quantitative preparation and a fairly transactional relationship with
learning. Get the answer and move on. My job, as I understand it, is to complicate that
relationship in a useful way. I want students to leave my courses not just with content
knowledge, but with something more durable: the ability to sit with a problem they cannot
immediately solve and keep working.

That sounds simple. It is not. Students who have been rewarded throughout their education
for speed and correctness do not easily tolerate confusion. Part of teaching is convincing
them that confusion is not a sign that something has gone wrong. It is often a sign that
something is starting to go right.

\stsection{II. What Foundational Means to Me}

I have spent most of my career teaching what the academy calls foundational or general
education courses. I do not think of these as lesser than upper-division work. I think of
them as load-bearing. The habits of mind that students develop in introductory physics, or
in a first-year science course, shape how they approach every technical problem they
encounter afterward.

This is also why I think the design of introductory courses matters so much and why so many
of them fall short. We tend to build them around coverage, a model driven more by tradition
and textbook architecture than by educational purpose. Faculty teach what they love: the
elegance of a derivation, the subtle mechanism connecting cause to effect. But those are
often the wrong things to lead with. Students do not yet have the scaffolding to see the
beauty through the complexity. What they encounter first \textit{is} the complexity, and many
of them conclude that the field simply is not for them. That conclusion is wrong, and the
course design helped them reach it.

The question I keep returning to is not how many chapters can be covered but what ideas
students will carry with them five years from now. Those outcomes do not require simplifying
the discipline. They require clarifying it.

\stsection{III. Active Learning in Practice}

I structure my courses around active engagement because I do not believe that students learn
physics by watching someone else do physics. Lecture has its place and I use it, but the core
of my courses is students working through problems, data, and laboratory tasks that require
interpretation and judgment rather than just measurement.

In practice, this means I spend a significant amount of class time circulating rather than
presenting. I ask questions more than I answer them. I have learned to be comfortable with
silence after a question, because the silence often means students are actually thinking. I
try to pose problems that require students to make a decision before they can proceed,
because decision-making is where understanding either holds up or breaks down.

I also try to make the reasoning visible. When I do work through a problem at the board, I
narrate my thinking, including the false starts, because I want students to see that
problem-solving is a process of iteration rather than a clean linear path from question to
answer. Expertise looks effortless from the outside. That appearance is misleading, and it
discourages students more than almost anything else we inadvertently do.

\stsection{IV. Productive Struggle}

The phrase I come back to most often, in my teaching and in my work leading the School of
Foundational Studies, is productive struggle. I use it deliberately, because not all struggle
is productive. Struggle that is too far beyond a student's current capacity is just
discouraging. Struggle calibrated correctly, that pushes students just past what they can do
comfortably, is where real learning happens.

My job is to design that calibration. It means building assignments that have genuine teeth
without being arbitrary. It means giving feedback that identifies the error in a student's
reasoning without simply supplying the corrected version. It means resisting the impulse to
rescue students too quickly when they get stuck, because the sticking point is often exactly
where the learning is located.

I am also honest with students about why I do this. The difficulty is not the obstacle; it is
the point. In an era when correct answers are increasingly easy to obtain, the value of an
engineering education is not the answers. It is the capacity to generate and evaluate them. I
try to make sure my courses are actually building that capacity rather than just assessing it
after the fact.

Introductory courses also carry more identity weight than we usually acknowledge. Students
use them to determine what kind of thinker they are and whether they belong in higher
education at all. When struggle is interpreted as evidence of incapability rather than as a
natural feature of learning, students do not just disengage from the material; they disengage
from the discipline. Designing courses that send a different message is not just good
pedagogy. It is a matter of access.

\stsection{V. Assessment and Feedback}

I try to make assessment a learning tool rather than just a measurement tool. This means
building in opportunities for revision, asking students to explain their reasoning rather than
just report their results, and designing assessments where partial credit is meaningful. A
student who can identify where their reasoning broke down has learned something. A student who
simply receives a score and moves on has not.

My feedback tends to be direct. I identify the conceptual error, not just the arithmetic one,
and I ask students to come to office hours when written feedback alone is not enough. I try to
make clear that I am interested in their development, not just in recording their performance.

\stsection{VI. Technology and AI in the Classroom}

My approach to technology, including AI, follows the same logic as the rest of my teaching.
The question is not whether a tool is available but whether its use in a given moment serves
the student's development or substitutes for it. A student who uses a calculator to skip
arithmetic they have already mastered is using a tool sensibly. A student who uses an AI
assistant to bypass reasoning they have not yet developed is shortchanging themselves, even if
the output looks correct.

I try to make that distinction explicit with students rather than managing it primarily
through policy. Rules about what tools are permitted matter, but they are not sufficient. The
more important work is helping students develop the judgment to know when a tool is helping
them learn and when it is helping them avoid learning. That judgment is itself a foundational
skill, and one that will only become more important.

\stsection{VII. What I Am Still Working On}

I do not have this figured out. I teach students with a wide range of preparation, a wide
range of reasons for being in my courses, and a wide range of responses to the kind of
challenge I try to create. What works for one group does not always work for another. I am
still learning how to calibrate difficulty across that range, how to build genuine equity into
active learning environments that can inadvertently reward students who are already
comfortable with ambiguity, and how to give students enough support that the struggle stays
productive rather than becoming discouraging.

I also think about the tension between coverage and depth. Physics curricula are dense, and
the pressure to cover material is real. I have made peace, mostly, with the idea that deep
engagement with fewer things is more valuable than shallow exposure to many, but I do not
always get the balance right.

I keep teaching because I find the problem genuinely interesting. Getting a student to see
something they could not see before, or to persist through a problem they wanted to abandon,
remains one of the more satisfying things I do. I expect my thinking about how to do that well
to keep evolving.

\signatureblock{Fall 2026}

\end{document}
